\subsection{Лекция 2 (12.09.26) МКР для решения уравнения Пуассона}
В тесте \cvar{poisson1-fdm} из файла \ename{poisson_fdm_test.cpp}
реализовано решение одномерного уравнения Пуассона с граничными условиями первого рода
(смотри теорию в \secref{sec:poisson_fdm}, объяснение примера в \secref{sec:poisson1d_prog}).
Проводится расчёт на сгущающихся сетках с количеством ячеек от 10 до 1000
и расчитываются среднеквадратичные нормы отклонения полученного численного решения от точного.
Решения сохраняются в vtk-файлы \ename{poisson1_fdm_n={}.vtk}.
Отталкиваясь от этой реализации необходимо:
\begin{enumerate}
\item Нарисовать график, на котором показать второй порядок сходимости схемы. При этом 
      увеличить диапазон рассматриваемых разбиений до того $N$, при котором поведение нормы сохраняет второй порядок сходимости.
      Проиллюстрировать численное решение для разных разбиений в сравнении с точным решением в Paraview.
\item Вместо условия первого рода использовать условие Неймана на левой границе расчётной области.
      Сравнить подходы первого \cref{eq:poisson1d_bc2_o1} и второго \cref{eq:poisson1d_bc2_o2} порядков,
      для каждого из которых нарисовать график сходимости.
\item Написать аналогичный тест, но для двумерной задачи Пуассона с граничными условиями первого рода.
      Проиллюстрировать решение в Paraview и показать второй порядок сходимости.
\end{enumerate}


\paragraph{Рекомендации}
\begin{itemize}
\item 
Свои программы следует оформлять в виде отдельных тестов (вместо того, чтобы модифицировать существующие).
Желательно, после того как программа заработает, сразу оставить несколько базовых \cvar{CHECK} проверок и
сделать локальный коммит, чтобы впоследствии легко распознавать и исправлять внесённые в дальнейшей работе ошибки.
К тому же наличие готовых тестов значительно облегчает рефакторинг кода.

При написании новых тестов следует переиспользовать уже написанный код, избегая копирования.
Для этого необходимо оформлять повторяющийся код в виде отдельных процедур и пользоваться механизмами наследования классов.

\item
Для иллюстрации одномерного решения в Paraview для п.1 см.~\secref{sec:paraview-1d}.

\item
Для программирования условий второго рода (п.2) вносить изменения нужно только в процедуры \cvar{Problem::approximate_lhs}, 
\cvar{Problem::approximate_rhs}. Вычислять граничное значение $q_a$ следует используя имеющееся точное решение.

\item
Для выполнения п.3 необходимо написать новый тест, со своим классом-решателем \cvar{Problem2D}. Для
уменьшения повторов в коде можно воспользоваться механизмом наследования классов.
В этом классе нужно определить двумерное точное решение \cvar{exact_solution} и расчитать соответсвующюю ему правую часть \cvar{rhs}.

\item
Двумерная задача Пуассона в частных производных будет иметь вид
\begin{equation*}
   -\left(\dfrq{u}{x} + \dfrq{u}{y}\right) = f(x, y).
\end{equation*}

Разностную схему для него удобно писать и использованием парного индекса
\begin{equation*}
   \frac{-u_{k[i-1,j]} + 2 u_{k[i, j]} - u_{k[i+1,j]}}{h_x^2} +
   \frac{-u_{k[i,j-1]} + 2 u_{k[i, j]} - u_{k[i,j+1]}}{h_y^2} =
   f_{k[i,j]},
\end{equation*}
где
\begin{equation}
    \label{eq:tasks_ij2k}
    k[i, j] = i + (n_x+1) j
\end{equation}
-- функция, переводящая парный $(i,j)$ индекс узла регурярной сетки ($i$) для оси x, $j$ для оси y) в сквозной индекс $k$
сеточного вектора, $n_x$ -- количество ячеек сетки в направлении x.

Функции, переводящие сквозной индекс в пару $i,j$, имеют вид
\begin{equation}
    \label{eq:tasks_k2ij}
    \begin{array}{ll}
        i[k] = {\rm mod}\left(k, \left(n_x+1\right)\right), & \text{// остаток от деления},\\
        j[k] = \lfloor k / \left(n_x+1\right)\rfloor,       &\text{// целая часть от деления}.\\
    \end{array}
\end{equation}

\item
Использовать класс \cvar{cfd::RegularGrid2D} для задания сетки.
Функции перевода индексов узлов из сквозных в парные и обратно реализованы в классе двумерной регулярной сетки:
\begin{itemize}
\item \cvar{cfd::RegularGrid2D::to_split_point_index}
\item \cvar{cfd::RegularGrid2D::to_linear_point_index}
\end{itemize}

\item
При вычислении весов $w_k$ для вычисления среднеквадратичного отклонения учесть
наличие граничных и угловых точек:
\begin{equation*}
    w_k = \begin{cases}
            h_x h_y / 4,  &\quad  \text{для угловых точек};\\
            h_x h_y / 2,  &\quad  \text{для граничных неугловых точек};\\
            h_x h_y,      &\quad  \text{для внутренних точек}.
    \end{cases}
\end{equation*}
Четрые угловые точки определяются как
\begin{equation*}
    i[k], j[k] = (0, 0), \, (0, n_y), \, (n_x, n_y), \, (n_x, 0)
\end{equation*}
Граничные неугловые точки:
\begin{equation*}
    \begin{array}{lll}
        i[k], j[k] =& \overline{1,n_x-1}, 0;    & \text{нижняя сторона},\\
                    & n_x, \overline{1,n_y-1};  & \text{правая сторона},\\
                    & \overline{1, n_x-1}, n_y; & \text{верхняя сторона},\\
                    & 0, \overline{1,n_y-1};    & \text{левая сторона}.
    \end{array}
\end{equation*}

\item
В случае, если решатель системы линейных уравнений не решает построенную матрицу, использовать функцию
\cvar{cfd::dbg::print} для отлаточной печати матрицы в консоль (размерность задачи должна быть небольшой).

\item
Для иллюстрации двумерного решения в Paraview использовать изолинии (\ref{sec:paraview-isolines}) и трёхмерные поверхности (\ref{sec:paraview-2d}).

\item
Следует иметь ввиду, что на графике сходимости
по оси абсцисс отложено линейное разбиение, вычисляемое как
$n=\sfrac{1}{h}$, где $h$ -- это характерный линейный размер ячейки.
Для двумерных сеток этот линейный размер
сетки можно вычислить через среднюю площадь ячейки $A$ как $h = \sqrt{A}$,
которую в свою очередь можно получить, разделив общую площадь на количество
ячеек: $A = \sfrac{|D|}{N}$. Тогда, в случае единичного квадрата, линейное разбиение будет равно $n = \sqrt{N}$.

\end{itemize}
